
\uuid{3p7H}
\titre{ Matrice dépendant d'un paramètre }

\niveau{L2}
\module{Algèbre}
\cpge{mpsi-mp2i.calcul-matriciel-systemes.a, pcsi.calcul-matriciel-systemes-lineaires.a, ptsi.calcul-matriciel-systemes-lineaires.a, mpsi-mp2i.polynomes-fractions-rationnelles.b, pcsi.polynomes.b, ptsi.polynomes.b, mpsi-mp2i.polynomes-fractions-rationnelles.c, pcsi.polynomes.c, ptsi.polynomes.c, bcpst1.matrices, bcpst1.polynomes-reels.a, mp-mpi.reduction-endomorphismes.c, pc.reduction.b, psi.reduction-endomorphismes.b, pt.reduction-endomorphismes.b, mp-mpi.reduction-endomorphismes.g, pc.complements-espaces-vectoriels.d, psi.complements-espaces-vectoriels.d, bcpst2.valeurs-propres-vecteurs-propres.a, bcpst2.valeurs-propres-vecteurs-propres.b}
\chapitre{Réduction d'endomorphisme}   			%Continuité, Groupes, Fonctions de plusieurs variables...
\sousChapitre{Diagonalisation}			%Optimisation, Diagonalisation d'une matrice, Calcul de dérivées partielles...

\theme{}				%Fonction de répartition, Division euclidienne de polynômes, ...
\auteur{Quentin Liard}
\datecreate{ {2026-03-26}}
\organisation{AMSCC}			%AMSCC, Exo7, ...
\difficulte{}			%1, 2, 3, 4 ou 5
\contenu{

\texte{
On considère la matrice

\[
A_\alpha=
\begin{pmatrix}
1 & 1 & 0\\
0 & 1 & 1\\
0 & 0 & \alpha
\end{pmatrix}.
\]
}

\begin{enumerate}

\item \question{Calculer le polynôme caractéristique de $A_\alpha$.}

\item \question{Déterminer le spectre de $A_\alpha$.}

\item \question{Pour quelles valeurs de $\alpha$ la matrice est-elle diagonalisable ?}

\item \question{Déterminer le polynôme minimal de $A_\alpha$ selon les valeurs de $\alpha$.}

\end{enumerate}

}